%%%
%
% $Autor: Wings $
% $Datum: 2021-05-14 $
% $Pfad: GitLab/MLEdgeComputer $
% $Dateiname: 
% $Version: 4620 $
%
% !TeX spellcheck = de_DE/GB
% !TeX program = pdflatex
% !BIB program = biber/bibtex
% !TeX encoding = utf8
%
%%%






\chapter{TinyML and Edge Computing}

\section{TinyML} 
``TinyML( Tiny Machine Learning) is a technology that can be used to develop embedded low power consuming devices to run both machine and deep learning models'' \cite{Gupta:2020c}. It requires simple procedures for training the model created. The amount of data required for creating the model is significantly less compared to other platforms. Tiny ML also works on battery powered devices and can be used for various applications including environmental monitoring. 

TinyML focuses on application design, development, deployment of different models which can be transformed in to a large scale ML model. This approach mainly removes the barriers of traditional style ML models and allows users to develop, deploy and implement applications seamlessly.” TinyML permits a wide range of new applications that traditional ML cannot deliver because of bandwidth, latency, economics, reliability, and privacy (BLERP) limitations.“ \cite{Reddi:2021}  
Some of the common applications of TinyML includes keyword spotting, wake word detection, person detection, gesture detection and much more. Keyword spotting generally refers to identification of words that typically act as part of a cascade architecture to kick-start or control a system, such as a mobile phone responding to voice commands \cite{Reddi:2021}. Visual wake words involve parsing image data to find an individual (human or animal) or object which can be used for security systems and for intelligent lighting. 


\section{Real Time} 

Real Time can be considered as a way of training the model by making it run through live data constantly in order to improve the model. This contradicts the traditional way of machine learning when machine learning engineers were dependent already existent data inputs for creating the model. 
A method by which we can implement real time learning in to a machine learning model is by continuously feeding it with a data stream and improving the model over time. This is achieved by using an event driven architecture. Event driven architecture is a modern design and development approach that centers about the events occurring in the model. Event-driven architectures create, detect, consume, and react to events.\cite{Hazelcast:2022}. By using a scalable event driven architecture, large number of events in real time can be responded, to which are suitable for loosely coupled softwares(For eg: web services). They also work well with unpredictable and non linear events by making it versatile. \cite{Hazelcast:2022}

\section{Edge computing}

The placement of computer and storage resources at the point where data produced is known as edge computing. This computes and stores  the data source at the network edge, which is optimal. In other words, instead of sending raw data to a main data centre for processing and analysis, this work is done right where the data is generated.
Edge computing is used to handle discrete tasks like determining if someone answered "yes" and responding appropriately. Instead of comparing the data with that on the web, the audio analysis is done on the edge. This drastically decreases expenses and complexity while also reducing the risk of data breaches.\cite{Bigelow:2021}

Edge computing has gained traction as a viable solution to a variety of issues related to transporting the massive amounts of data that today's businesses generate and consume. It's not just a matter of quantity, it's also a matter of time; applications increasingly rely on processing and responses that are time-sensitive.
